\subsection{Групповой полет (только на дронах)}
\label{sec:python-swarm}

Управление роем (Swarm) — сложная, но захватывающая задача. Главное здесь — синхронизация и единая система координат.

\subsubsection*{Архитектура роя}
Мы будем использовать централизованную архитектуру:
\begin{itemize}
    \item \textbf{Один компьютер (Наземная станция)} управляет всеми дронами.
    \item Каждый дрон имеет свой уникальный IP-адрес.
    \item Программа создает отдельный поток (\texttt{Thread}) для каждого дрона.
\end{itemize}

\begin{infobox}[Что с роем в симуляторе]
В симуляторе каждый новый объект \verb|Pioneer| с другим именем или адресом привязывается к следующему свободному дрону сцены, поэтому логику роя можно попробовать на сцене с несколькими дронами. Однако примеры ниже рассчитаны на настоящие IP-адреса и потоки \verb|threading|, а в браузере потоки выполняются по очереди и общий момент команд не гарантируют; в симуляторе рой не проверялся, и результат надёжно получается только на настоящих дронах.
\end{infobox}

\subsubsection*{Пример: Синхронный взлет}
В файле \texttt{group\_flight\_examples/swarm\_OPT.py} показан пример, где несколько дронов выполняют задачи параллельно.

\textbf{Ключевые моменты:}
1. Создаем список IP-адресов.
2. Для каждого IP создаем объект \texttt{Pioneer}. Адрес обязательно передается по имени --- \texttt{Pioneer(ip=ip)}: первый позиционный параметр конструктора называется \texttt{name}, поэтому запись \texttt{Pioneer(ip)} задаст дрону имя, а подключение пойдет на адрес по умолчанию.
3. Используем \texttt{threading.Thread} для запуска функций управления одновременно.

\begin{codebox}[Псевдокод роя]
\begin{verbatim}
drones = [Pioneer(ip=ip1), Pioneer(ip=ip2)]

def mission(drone):
    drone.arm()
    time.sleep(1)   # пауза между arm и takeoff обязательна
    drone.takeoff()
    time.sleep(5)
    drone.land()

# Запуск потоков
threads = []
for d in drones:
    t = threading.Thread(target=mission, args=(d,))
    t.start()
    threads.append(t)

# Ожидание завершения
for t in threads:
    t.join()
\end{verbatim}
\end{codebox}

\begin{warnbox}[Безопасность роя]
При групповых полетах риск столкновения возрастает многократно!
1. Всегда проверяйте координаты старта.
2. Расстояние между дронами должно быть не менее 1.5 метров.
3. Держите палец на кнопке аварийного отключения (или имейте общий скрипт \texttt{disarm\_all}).
\end{warnbox}

\begin{taskbox}[Танец]
Только на настоящих дронах (на сцене симулятора --- по желанию, см. врезку выше). Напишите программу для двух дронов с разными IP: дрон А взлетает; дрон Б ждёт 5~секунд, затем взлетает; дрон А садится; дрон Б садится. Каждым дроном управляет своя функция в своём потоке, а моменты событий согласуйте через \verb|threading.Event| (А ставит флаг «взлетел», Б ждёт его). Расстояние между точками старта --- не меньше 1.5~м.
\expected{дроны взлетают с интервалом около 5~с и садятся в порядке А, Б; если один дрон не подключился, программа сообщает об этом и не поднимает второй.}
\end{taskbox}

\subsection*{Итоги части о зрении и рое}
\begin{itemize}
    \item Кадр может не прийти (\verb|None|), метки в кадре может не быть (\verb|ids is None|) --- это разные случаи, и оба проверяют до любых вычислений.
    \item ID метки называет рисунок, а не место: смысл ID задаёт таблица в программе, положение метки --- только кадр.
    \item Ошибка метки относительно центра кадра --- измерение в пикселях; из неё скорость делает регулятор с коэффициентом, пределом, мёртвой зоной и обработкой потери метки.
    \item Расстояние в метрах даёт \verb|solvePnP| по размеру метки и параметрам камеры; размер и порядок углов должны совпадать с реальными.
    \item В симуляторе работают камера, ArUco, наведение и \verb|solvePnP| с идеальной камерой; калибровка, \verb|imwrite|, MediaPipe и настоящий рой --- только на компьютере и дроне.
\end{itemize}

\begin{questionbox}[Контрольные вопросы]
\begin{enumerate}
    \item Чем отличается \verb|frame is None| от \verb|ids is None| и что делает программа в каждом случае?
    \item Метка правее и ниже центра кадра. Какие знаки у $e_x$ и $e_y$ и почему у нижней камеры симулятора $v_y$ берут с противоположным знаком к $e_y$?
    \item Почему нельзя двигаться по ошибке, измеренной несколько кадров назад?
    \item Что изменится в расстоянии из \verb|tvec|, если в коде размер метки указан вдвое меньше настоящего?
    \item Зачем в групповом полёте проверять подключение всех дронов до взлёта первого?
\end{enumerate}
\end{questionbox}
